\documentclass{article}
\usepackage{amsmath,amssymb,hyperref}
\title{How much identity does a two-scan forecast need?}
\date{}
\begin{document}
\maketitle
A historical scan archive can retain each reading exactly once, yet
lose the association needed to decide where to stage a reserve.
This note extends the two-scan identification model in
\href{https://github.com/eternal-november/eos-research-programmes/pull/398}{PR~398}.
The proposed intervention is a coarse physical tag at the two scan cuts,
not a unique identity for each run. The guarantee is exact for the two
calibrated rosters below, not a statistical estimate of their prevalence.

\paragraph{Information and mechanism.}
The two possible persistent source modes call for opposite placements.
Each has a fixed clean roster of paired binary scans. The controller
knows the two candidate clean tables, not the true mode or pairs. Each
mode agrees on the two binary margins within every tagged block. For a
fixed equipment setting, the controller receives only a table per
reported block, or the single pooled table if tags are absent. True
identifiers and exact timestamps are not available to the policy.
A join may permute right readings within a block. A single total bound
$K$ limits displaced right identities across the whole roster; the
positions can depend on the source. Readings are neither changed, dropped
nor duplicated. In the untagged control all within-roster pairings are
allowed. In the tagged arm the authentic epoch marker attached by both
read stations agrees for every genuine pair, so no admissible join
crosses it. The tag must arrive before historical data are used to stage
a reserve. Were a tag unreliable, a different ambiguity set would apply.

\paragraph{Exact reachability.}
In one block let $a,b,c,d$ be the true counts at $00,01,10,11$. With
fixed margins a candidate linked table is determined by its $00$ cell
$x$. The least number of displaced identities is $2|x-a|$ when
$a-\min(a,d)\le x\le a+\min(b,c)$, and the table is impossible outside
this interval. To increase $a$, swap right identities between cells
$01$ and $10$, one disjoint swap per unit. To decrease $a$, use $00$
and $11$. Conversely an increase or decrease by one requires one
changed right value in each of the two left strata; a displacement
of an identity cannot change more than one value. Thus for blocks
the exact joint uncertainty set consists of the common fixed-margin
box intersected with the integer ball
\[
      \sum_j |x_j-a_j|\le \lfloor K/2\rfloor .
\]
A single pooled error interval or one full allowance per block is not
that set. Swaps at two unrelated places still spend the same global
budget. The result does not cover overlapping clock windows that
allow only some specific pairs within or across blocks.

\paragraph{The forecast boundary.}
For two clean center vectors $a^{(0)},a^{(1)}$, their possible
observations are disjoint if and only if
\[
  D=\sum_j |a_j^{(0)}-a_j^{(1)}|>2\lfloor K/2\rfloor .
\]
Necessity for common observations follows from the triangle inequality.
Sufficiency follows by walking along an integer coordinate path between
the two centers. Each point stays in the fixed-margin box; a midpoint
attains the two radius bounds when $D\le 2\lfloor K/2\rfloor$.
Nearest center recovers the right placement whenever the sets are
disjoint. If they meet, some common observation forces one of the two
modes to have error at least $1/2$ for a randomized placement, and
fair guessing attains that bound. This is worst-case error over
the two fixed rosters and all allowed joins, not a posterior risk.

\paragraph{A tag can help without fixing a single ambiguous pair.}
There are two batches of 80 runs each, with balanced scans in both.
Mode 0 has $00$ counts $(40,16)$; mode 1 has $(0,24)$. Each correctly
linked pooled roster gives 30 percent scan-parity error relative to
its true mode. The pooled center distance is 32, since the two batch
differences partially cancel. A one-to-one join moving 32 records
can make both pooled tables $(40,40,40,40)$; no edited reading is
needed. With the tags, the center distance is 48. The same budget 32
cannot erase the forecast even though both blocks still admit arbitrary
confusion internally. This remains true through $K=47$. At $K=48$,
disjoint swaps in each world realize the common vector $(20,20)$.

A concrete comparison charges one optional maintenance window to
operate the epoch marker across both stations. In cost units where
missing the right reserve placement costs 10, the best pooled
guarantee at $K=32$ costs 5 in worst-case expected placement loss
without that work. If collecting the reliable tags costs 1.2 and
does not alter the future fault law or available placements, the
tagged guarantee costs 1.2. If it costs more than 5, the experiment
is not worth that loss reduction. At $K=48$ these two tags cannot
justify a positive acquisition cost on the worst-case criterion.
The tags and costs here are stipulated model controls, not deployed
hardware or measured values. A secure clock or authenticated
occurrence ID would need its own physical validation before use.

\paragraph{Comparison and scope.}
Ordinary one-to-one record linkage and sensitivity analysis already
model such errors; see
\href{https://doi.org/10.1093/ije/dyx177}{Harron et al.} and
\href{https://www.sciencedirect.com/science/article/pii/S0378375821001063}{Scholtus et al.}.
The finite same-information permutation and feasibility check agrees
with the formula. A tag may change the data available, and that is the
comparison here, not a gain over a planner given the same data.
The untagged observation includes neither a join-confidence score nor
identity metadata; adding truthful ancillary information can break the
lower bound. A clean tag supplies a smaller candidate set, not a claim
of unique reconstruction or general R3b closure.
\end{document}
